\chapter*{Abstract}
\addcontentsline{toc}{chapter}{ABSTRACT}

This proposal describes BiasAperture, a diagnostic framework for evaluating demographic disparities in facial-analysis model predictions and producing traceable research evidence. It specifies five modules for data ingestion, model interfacing, fairness metrics, surrogate attribution, and report generation, and discusses four metrics: Demographic Parity Difference, Equalized Odds Difference, Equal Opportunity Difference, and Disparate Impact Ratio. Fairlearn and AIF360 are proposed as separate metric implementations; their use does not constitute independent assessment. The reviewed implementation computes subgroup confidence intervals with an empirical percentile bootstrap that defaults to $B=200$; adapters do not pass a larger resample budget. This does not meet the proposed NFR-002 minimum of 1{,}000 resamples. A separate global-summary path uses a BCa bootstrap routine, which does not close the subgroup-interval gap. The demographic logistic surrogate summarizes associations between demographic indicators and model predictions; it does not explain image pixels, identify causes or confounders, or decompose a disparity. UTKFace was excluded from the completed scope; FairFace is the retained benchmark. The proposed report crosswalk is a research aid, not evidence of regulatory compliance or a compliance determination. BiasAperture is diagnostic in scope and does not modify model weights, apply mitigations, or generate demographic images.

\vspace{1cm}
\noindent
\textbf{Keywords:} Algorithmic Fairness, Bias Auditing, Demographic Parity, Disparate Impact, EU AI Act, Facial Analysis, Fairness Metrics, NIST AI RMF, Surrogate Attribution
